\section{\texorpdfstring{THE UNIFORMITY OF \(\mathfrak{S}\)-CONVERGENCE}{THE UNIFORMITY OF S-CONVERGENCE}}

Notation. If X and Y are any two sets, we recall that the set of all mappings of X into Y is denoted by \(\mathcal{F}(\mathbf{X};\mathbf{Y})\) , and may be identified with the product set \(\mathbf{Y}^{\mathbf{x}}\) (Set Theory, Chapter II, § 5, no. 2). For each subset H of \(\mathcal{F}(\mathbf{X};\mathbf{Y})\) and each \(x\in \mathbf{X}\) , we shall denote by \(\mathrm{H}(x)\) the set of elements \(u(x)\in \mathbf{Y}\) as \(u\) runs through H. If \(\Phi\) is a filter base on \(\mathcal{F}(\mathbf{X};\mathbf{Y})\) , we denote by \(\Phi (x)\) the filter base on Y formed by the sets \(\mathrm{H}(x)\) as H runs through \(\Phi\) . Finally, we recall that, for each \(u\in \mathcal{F}(\mathbf{X};\mathbf{Y})\) and each subset A of X, \(u|\mathrm{A}\) denotes the restriction of \(u\) to A, which is a mapping of A into Y; if H is a subset of \(\mathcal{F}(\mathbf{X};\mathbf{Y})\) , H\textbar A will denote the set of restrictions \(u|\mathrm{A}\) of functions \(u\in \mathrm{H}\) .

